% Full limiting-pattern proof, including both orientations and a comparison.

\section{Kamp's limiting past and future lemma}
\label{sec:annex-p09-kamp-iii2-lemma2}

\begin{annexitem}{annex:p9.kamp.iii2.lemma2}{thm:p09-kamp-iii2-lemma2}{Strict limits and the complete finite-chain splice}
Use the language, named-variable expressibility, and basic word realizations
already fixed in the body and in \cref{annex:p9.kamp.iii2.lemma1}.
The proof below develops the two finite-chain decompositions used in the
source's Lemma~2 and supplies the future direction explicitly
\cite[ch.~III, pp.~57--59, 63--67, 70--73]{Kamp1968}.

\begin{proof}
Fix an arbitrary structure
\(\mathcal M=\langle T,<,(P_i^{\mathcal M})_i\rangle\) over a nonempty
strict linear order, and let \(t\in T\) be the current point.
Every temporary variable introduced below is fresh from its ambient formula.
The \(D\)-translation identifies truth of a temporal formula at a point
with truth of its named-variable first-order image there.

\paragraph{The two derived limit operations.}
The formula \(\neg(\eta\land\neg\eta)\) is true at every point.
Expanding \(P'(\eta)=\neg H'(\neg\eta)\) through the strict-since
semantics therefore gives
\[
\begin{split}
 \mathcal M,t\models P'(\eta)
 &\quad\Longleftrightarrow\quad
 \neg\exists a<t\,\forall z\,
       (a<z<t\ \Rightarrow\ \mathcal M,z\models\neg\eta)\\
 &\quad\Longleftrightarrow\quad
 \forall a<t\,\exists z\,
       (a<z<t\ \text{ and }\ \mathcal M,z\models\eta).
\end{split}
\tag{1}
\]
The final equivalence is ordinary classical quantifier negation, including
an empty open interval.  A witness \(a\) with no point between it and
\(t\) satisfies the intervening universal condition and makes \(P'\)
false.  If there is no \(a<t\), both the final universal condition and
\(P'\) are true.

Expanding \(F'(\eta)=\neg G'(\neg\eta)\) through strict until gives
\[
\begin{split}
 \mathcal M,t\models F'(\eta)
 &\quad\Longleftrightarrow\quad
 \neg\exists b>t\,\forall z\,
       (t<z<b\ \Rightarrow\ \mathcal M,z\models\neg\eta)\\
 &\quad\Longleftrightarrow\quad
 \forall b>t\,\exists z\,
       (t<z<b\ \text{ and }\ \mathcal M,z\models\eta).
\end{split}
\tag{2}
\]
Here also empty intervals and a missing future domain are included.
If \(\eta\) expresses \(\theta(w)\), substitute its truth equivalence
at each quantified point in (1) or (2).  The results are exactly the two
formulas \(\partial^-\theta(w)\) and \(\partial^+\theta(w)\), with
\(w\) assigned \(t\).  This proves the body's claim about the formula
operations, without a density or endpoint assumption.

\paragraph{The exact chain conditions of the limiting classes.}
Let \(\sigma=(A_0,B_1,A_1,\ldots,B_n,A_n)\).  For this proof, write
\(C^-_\sigma(t)\) for the following condition: for every \(a<t\),
there is a realization \(z_0<\cdots<z_n\) of \(\sigma\) such that
\(a<z_0\) and \(z_n<t\).  Write \(C^+_\sigma(t)\) for the condition
that for every \(b>t\) such a realization has \(t<z_0\) and \(z_n<b\).
For \(n=0\), a realization consists of one point satisfying \(A_0\).
These are finite first-order conditions when \(\sigma\) is fixed:
quantify the \(n+1\) points, their strict order, each endpoint predicate,
and each universal open-interval predicate.

The basic-pattern realization equivalence in
\cref{annex:p9.kamp.iii2.lemma1} shows that a formula of
\(\operatorname{Lp}_n(w)\) with word \(\sigma\) is true at \(t\)
exactly when \(C^-_\sigma(t)\) holds.  Indeed, for \(n\geq1\) its
two witnesses satisfy \(a<u<v<t\).  The constituent basic pattern is
equivalent, at these ordered endpoints, to a realization beginning at
\(u\) and ending at \(v\).  Every intermediate point also lies in
\((a,t)\).  Conversely, a realization wholly in that interval supplies
\(u=z_0,v=z_n\); they are distinct and ordered because \(n\geq1\),
and the basic-pattern equivalence verifies the constituent formula.
At \(n=0\) the claim is the definition of \(\partial^-A_0\).

The future case has the same endpoint verification with its actual bounds:
its witnesses satisfy \(t<u<v<b\).  A realization of its constituent
lies wholly inside \((t,b)\); conversely a realization inside \((t,b)\)
supplies the first and last endpoints and verifies that constituent.  At
stage zero it is exactly \(\partial^+A_0\).  Thus the future formulas
are characterized by \(C^+_\sigma(t)\).

The same basic lemma identifies \(E^+_\sigma\) at a point with the
existence of a realization of the whole word beginning there, and identifies
\(E^-_\sigma\) with the existence of one ending there.  The chosen binary
past or future association does not change this assertion: the guarded
same-word equivalence has already been proved.  Repeated predicate letters
remain repeated positions of one word and require no independent valuations.

\paragraph{Past decomposition: both directions and every crossing.}
For \(n\geq1\) we claim
\[
 C^-_\sigma(t)\quad\Longleftrightarrow\quad
  (\mathcal M,t\models P'(E^+_\sigma))\ \text{ and }\
                                    C^-_{\sigma^-}(t).
\tag{3}
\]
Suppose first that \(C^-_\sigma(t)\) holds, and take any \(a<t\).
A realization \(z_0<\cdots<z_n\) inside \((a,t)\) has its first point
there, and that point satisfies \(E^+_\sigma\).  Equation (1), for every
such \(a\), proves the first conjunct.  Deleting \(z_0\) leaves a
realization of the suffix word \(\sigma^-\) inside the same interval.
At \(n=1\) this suffix is the one point \(z_1\) satisfying \(A_1\).
This proves the second conjunct in all cases.

Conversely, assume the two conjuncts on the right of (3), and take an
arbitrary \(a<t\).  By (1), choose \(u_0\) with \(a<u_0<t\)
satisfying \(E^+_\sigma\), and choose its whole-word realization
\[
                     u_0<u_1<\cdots<u_n.
\]
If \(u_n<t\), it is already the required realization inside \((a,t)\).
Otherwise \(u_n\geq t\).  Let \(k\) be the greatest index for which
\(u_k<t\).  This finite greatest index exists because \(u_0<t\), and
it satisfies
\[
             0\leq k<n,\qquad u_k<t\leq u_{k+1}.
\tag{4}
\]
In particular \(k=n-1\) is permitted.  Apply \(C^-_{\sigma^-}(t)\)
to the cut \(u_k<t\).  It supplies a suffix realization
\[
                    u_k<v_1<v_2<\cdots<v_n<t,
\]
where \(v_i\) satisfies \(A_i\), and the intervals between its
successive points satisfy \(B_i\) for \(2\leq i\leq n\).
Form the new sequence
\[
                  u_0,\ldots,u_k,v_{k+1},\ldots,v_n.
\tag{5}
\]
It has precisely \(n+1\) points.  It is strictly increasing, begins
above \(a\), and ends below \(t\).  Each position retains its required
predicate: old positions through \(k\) satisfy their \(A_i\), and
the suffix positions satisfy the remaining \(A_i\).

For every gap with index \(i\leq k\), the old realization supplies
\(B_i\).  For every gap with \(i\geq k+2\), the suffix realization
supplies \(B_i\).  The sole new gap is
\((u_k,v_{k+1})\).  By (4),
\[
       u_k<v_{k+1}<t\leq u_{k+1},\qquad
       (u_k,v_{k+1})\subseteq(u_k,u_{k+1}).
\]
The old gap therefore supplies exactly \(B_{k+1}\) throughout this new
gap.  This verifies every realization condition for (5).

When \(k=0\), the new sequence retains only the old first point and
replaces its whole remaining part.  When \(k=n-1\), it replaces only
the old last point.  When \(u_{k+1}=t\), the inclusion of open intervals
still holds and no predicate is demanded at \(t\).  These include all
boundary positions and equality with the current point.  At \(n=1\),
the argument uses a single suffix point and the one new gap, so it supplies
the source's special one-interval case as well.  Since \(a<t\) was
arbitrary, (5) proves \(C^-_\sigma(t)\), completing (3).

\paragraph{Future decomposition: the full counterpart.}
For \(n\geq1\) the future identity is
\[
 C^+_\sigma(t)\quad\Longleftrightarrow\quad
  (\mathcal M,t\models F'(E^-_\sigma))\ \text{ and }\
                                    C^+_{\sigma^+}(t).
\tag{6}
\]
For its forward direction, at any cut \(b>t\) choose a realization
\(z_0<\cdots<z_n\) inside \((t,b)\).  Its last point satisfies
\(E^-_\sigma\), so (2) proves the first conjunct.  Deleting that last
point leaves a realization of the prefix \(\sigma^+\) inside \((t,b)\).
For \(n=1\), this is the single point satisfying \(A_0\).  Hence the
second conjunct holds.

For the reverse direction assume both conjuncts and fix \(b>t\).
Equation (2) supplies a point \(u_n\) in \((t,b)\) satisfying
\(E^-_\sigma\), together with a realization
\(u_0<\cdots<u_n\) ending there.  If \(u_0>t\), all points already
lie inside \((t,b)\).  Otherwise \(u_0\leq t\).  Let \(k\) be the
least index whose point is greater than \(t\).  Then
\[
              1\leq k\leq n,\qquad u_{k-1}\leq t<u_k.
\tag{7}
\]
Apply \(C^+_{\sigma^+}(t)\) at the cut \(u_k>t\).  It gives a
prefix realization
\[
                    t<v_0<\cdots<v_{n-1}<u_k.
\]
Join its initial portion to the old final portion:
\[
                   v_0,\ldots,v_{k-1},u_k,\ldots,u_n.
\tag{8}
\]
This again has \(n+1\) points in the correct strict order, wholly between
\(t\) and \(b\).  Each point has its required \(A_i\) predicate.
The gaps with indices below \(k\) come from the new prefix and those
above \(k\) from the old realization.  For the new gap, (7) gives
\[
 u_{k-1}\leq t<v_{k-1}<u_k,\qquad
       (v_{k-1},u_k)\subseteq(u_{k-1},u_k).
\]
Thus the old gap supplies \(B_k\).  This verifies (8) as a whole-word
realization inside \((t,b)\).

The value \(k=1\) replaces only the old first point, and \(k=n\)
replaces all but the old last point.  Equality \(u_{k-1}=t\) causes no
change to the open-interval inclusion.  At \(n=1\) there is one new
prefix point and one gap.  These are all future cases, including the two
extreme crossing positions.  Since the cut \(b\) was arbitrary, (6)
follows.  Neither the past nor the future proof inserted an order point
without one of its stated existential hypotheses.

\paragraph{Induction giving the actual temporal formulas.}
We prove simultaneously that \(\Lambda^-_\sigma\) expresses
\(C^-_\sigma\) and \(\Lambda^+_\sigma\) expresses \(C^+_\sigma\),
by induction on the number \(n\) of intervals in the word.  At zero,
the two conditions are \(\partial^-A_0\) and \(\partial^+A_0\),
so (1)--(2) prove the assertion for the two declared seeds.

For \(n\geq1\), Lemma~1 already supplies the stated meanings of
\(E^+_\sigma\) and \(E^-_\sigma\).  The induction hypothesis supplies
the meanings of \(\Lambda^-_{\sigma^-}\) and
\(\Lambda^+_{\sigma^+}\), since their words have \(n-1\) intervals.
Equation (3) now identifies
\(P'(E^+_\sigma)\land\Lambda^-_{\sigma^-}\) with the full past
condition, and (6) identifies
\(F'(E^-_\sigma)\land\Lambda^+_{\sigma^+}\) with the full future
condition.  These are exactly the recurrences in the body.  The chain
characterization above transfers their meanings to each formula in its
corresponding limiting class.

Reading the finite occurrence word and applying these recurrences is an
explicit construction.  Each step deletes one interval, while its basic
formula is obtained by the already given finite since/until recurrence.
The derived limit operations also expand to finite \(\mathsf{TL}_1\)
formulas.  Only predicate letters of the word are used, including when a
letter occurs at several positions.  The \(D\)-translation, followed by
fresh bound-variable renaming and replacement of its free \(t_0\) by
\(w\), gives the required named-variable expressibility equivalence.

\paragraph{Endpoint and crossing calculations.}
At a first point there are no past cuts.  Thus every \(C^-_\sigma\)
holds, and (1) makes every past limit factor true.  At a last point there
are no future cuts, and every \(C^+_\sigma\) and every factor from (2)
is true.  If \(a\) is an immediate predecessor of \(t\), the cut
\(a<t\) has no point in its intervening interval.  It cannot contain
even a one-point realization, so all past limiting conditions fail.
Equation (1) gives the same failure for the temporal construction.  An
immediate successor gives the identical future failure through (2).
Consequently the singleton and adjacency claims in the body hold without
assumptions on the predicates.  Empty gaps between points of a realization
still satisfy their interval predicates vacuously; this does not create a
missing endpoint witness.

For the body's one-interval counterexample on \(\mathbb R\) at zero,
each cut \(a<0\) admits \(u=a/2\), where \(A(u)\) and
\(U(\widehat C,\widehat B)\) hold using endpoint 1.  Hence its first
past conjunct is true.  No point of \(C\) lies below zero, so the second
past conjunct and the entire past limiting pattern fail.  In the future
counterpart, for a cut \(b>0\) choose \(v=b/2\).  It is in \(C\)
and the chain starting at \(-1\) witnesses the since formula.  No point
of \(A\) lies above zero, so the second future conjunct and the future
limiting pattern fail.  The second conjuncts are therefore necessary.

With all predicates true on \(\mathbb R\), each cut admits any required
finite number of increasing points in its open interval, so both
decomposition premises hold.  In the displayed past crossing calculation,
\(n=2\), \(u_0=-3,u_1=-1,u_2=t=0\), and the last strictly past
index is \(k=1=n-1\).  The suffix \(v_1=-1/2,v_2=-1/4\) is above
\(u_1\); (5) is exactly \((-3,-1,-1/4)\).  Its new gap is contained
in \((-1,0)\).  This directly verifies the last crossing and equality case.
For a future first crossing take \((u_0,u_1,u_2)=(0,1,3)\), current
point zero and cut 4.  The prefix \((1/4,1/2)\) lies below \(u_1=1\),
and (8) gives \((1/4,1,3)\).  For a future last crossing take
\((-2,0,1)\) at zero with cut 2; the same prefix gives
\((1/4,1/2,1)\).  Both satisfy every endpoint and interval condition.

\paragraph{The current-point comparison.}
Consider \(\operatorname{CurF}_A(w)\) as separately
defined in the body.  If \(t\) has no future point, its universal formula
is true.  The temporal formula \(U(\tau_A,\tau_A)\) holds exactly
when a future point exists, since both occurrences of \(\tau_A\) are
tautologies.  Its negation therefore expresses this first case.

If a future point exists, instantiate this formula at one such cut.
Its matrix implies \(A(t)\).  At each future cut its remaining requirement
is that some point lie strictly between the cut and \(t\).  By (2), this
is exactly \(F'(\tau_A)\).  Conversely, \(A(t)\) together with
\(F'(\tau_A)\) supplies the witness and the unchanged current-point
predicate at every future cut.  This proves the displayed disjunction as an
expression for the current-point condition, including a last point.
On \(\mathbb R\) at zero with \(A=\{0\}\), choose \(z=x/2\)
for each \(x>0\): the current-point formula is true.  The source's base is
false, since there is no point of \(A\) in any strictly future interval.
This verifies the stated distinction directly.  This additional comparison
does not change the source's witness-point bases or the limiting classes
whose expressibility was proved above.

The entire argument is over \(\mathsf{Lin}\).  Completeness of the order,
the later elementary-class substitution theorem, and the remaining interval
lemmas have not been used.  Their required proofs, including the next source
Lemma~3, remain separate from the lemma proved here.
\end{proof}
\end{annexitem}
